\section{What attaches one layer to another}\label{x:extensions}

The additive invariant has reached its limit. For finite-length modules,
an exact sequence
\[
0\longrightarrow A\xrightarrow{i}E\xrightarrow{q}B\longrightarrow0
\]
always gives $[E]=[A]+[B]$, whether or not $E$ is a direct sum. To learn
more, we must retain the way $A$ sits in $E$ and the way $E$ maps onto $B$.
The maps are part of the object we now study.

Throughout this section $R$ is a ring with identity, modules are left
$R$-modules, and every map between modules is $R$-linear. Direct sums carry
componentwise operations. An \emph{extension of $B$ by $A$} means the whole
displayed short exact sequence, with $A$ and $B$ fixed. We will first
construct extensions, then calculate a family of them, and finally use
them to detect when a map cannot be lifted.

\subsection{Keeping the endpoints fixed}

Two extensions $(i,E,q)$ and $(i',E',q')$ are \emph{equivalent} if there is
a map $h:E\to E'$ making this diagram commute:
\[
\begin{tikzcd}[column sep=large]
0\arrow[r]&A\arrow[r,"i"]\arrow[d,equal]
 &E\arrow[r,"q"]\arrow[d,"h"]&B\arrow[r]\arrow[d,equal]&0\\
0\arrow[r]&A\arrow[r,"i'"]&E'\arrow[r,"q'"]&B\arrow[r]&0.
\end{tikzcd}
\]
Thus $hi=i'$ and $q'h=q$. No separate invertibility assumption on $h$ is
needed.

\begin{lemma}\label{x:marked-iso}
Every such $h$ is an isomorphism.
\end{lemma}
\begin{proof}
If $h(e)=0$, then $q(e)=q'h(e)=0$, so $e=i(a)$ for some $a\in A$.
Now $0=hi(a)=i'(a)$, and injectivity of $i'$ gives $a=0$ and $e=0$.
Hence $h$ is injective.

Given $e'\in E'$, choose $e\in E$ with $q(e)=q'(e')$. This is possible
because $q$ is surjective. The difference $e'-h(e)$ lies in $\Ker q'$, so
it is $i'(a)$ for some $a\in A$. Then $h(e+i(a))=e'$. Thus $h$ is
surjective. The inverse of a bijective linear map is linear: apply $h$ to
the proposed addition and scalar identities and use its injectivity.
\end{proof}

Identity maps, compositions, and inverses now show that equivalence is
indeed an equivalence relation. The word \emph{fixed} matters: an
isomorphism between the middle modules that changes an endpoint
identification need not give an equivalence of extensions.

The \emph{split extension} is
\[
0\longrightarrow A\xrightarrow{a\mapsto(a,0)}A\oplus B
 \xrightarrow{(a,b)\mapsto b}B\longrightarrow0.
\]
An extension is called split when it is equivalent to this one.

\begin{proposition}\label{x:splitting}
An extension is split if and only if there is a linear map $s:B\to E$
with $qs=\id_B$.
\end{proposition}
\begin{proof}
Such a map is a \emph{section} of $q$. Given a section, the map
\[
A\oplus B\longrightarrow E,\qquad (a,b)\longmapsto i(a)+s(b)
\]
commutes with both endpoint maps and is therefore an isomorphism by
Lemma~\ref{x:marked-iso}. Conversely, an equivalence with the split
extension transports its section $b\mapsto(0,b)$ to a section of $q$.
\end{proof}

A choice of a lift of each individual element of $B$ is insufficient.
The lifts must form a linear map. Failure of that compatibility is
precisely what a nonsplit extension records.

\subsection{Change the quotient: pull back}

Suppose we want to compare an extension with a map $f:P\to B$. We need
an object that simultaneously records an element of $E$ and an element
of $P$ which have the same image in $B$. Define
\[
f^*E=E\times_B P
 =\{(e,p)\in E\oplus P:q(e)=f(p)\}.
\]
This is a submodule: it is the kernel of the linear map
$(e,p)\mapsto q(e)-f(p)$. Put
\[
j(a)=(i(a),0),\qquad \pi(e,p)=p,\qquad u(e,p)=e.
\]
These maps form the diagram
\[
\begin{tikzcd}[column sep=large]
0\arrow[r]&A\arrow[r,"j"]\arrow[d,equal]
 &f^*E\arrow[r,"\pi"]\arrow[d,"u"]&P\arrow[r]\arrow[d,"f"]&0\\
0\arrow[r]&A\arrow[r,"i"]&E\arrow[r,"q"]&B\arrow[r]&0.
\end{tikzcd}
\]
The left square commutes because $uj=i$. The right one commutes because
membership in $f^*E$ says exactly that $qu=f\pi$.

The top row is exact. The map $j$ is injective because $i$ is. If
$\pi(e,p)=0$, then $p=0$ and $q(e)=0$, so $(e,p)=(i(a),0)=j(a)$
for a unique $a$. Conversely every $j(a)$ maps to zero. Finally, for
any $p\in P$, surjectivity of $q$ supplies $e$ with $q(e)=f(p)$, so
$\pi(e,p)=p$. No assumption that $f$ is injective or surjective was used.

The right square is a pullback in the precise universal sense. If maps
$r:T\to E$ and $s:T\to P$ satisfy $qr=fs$, the unique map to $f^*E$
with coordinates $r$ and $s$ is
\[
T\longrightarrow f^*E,\qquad t\longmapsto(r(t),s(t)).
\]
The compatibility puts its image in $f^*E$; coordinatewise operations
make it linear; its two coordinates force uniqueness.

\subsection{Change the subobject: push out}

Now suppose instead that $g:A\to A'$ is given. We want to replace the
embedded $A$ by $A'$ while preserving the quotient $B$. In $A'\oplus E$
we therefore identify $(g(a),0)$ with $(0,i(a))$. The required relation
submodule is
\[
L_g=\{(g(a),-i(a)):a\in A\},
\qquad
g_*E=(A'\oplus E)/L_g.
\]
Here $L_g$ is the image of a linear map, so it is a submodule and the
quotient is defined. Brackets below denote classes modulo $L_g$. Define
\[
i_g(a')=[a',0],\qquad u_g(e)=[0,e],\qquad q_g([a',e])=q(e).
\]
The last map is well-defined because $q(i(a))=0$. The diagram is
\[
\begin{tikzcd}[column sep=large]
0\arrow[r]&A\arrow[r,"i"]\arrow[d,"g"]
 &E\arrow[r,"q"]\arrow[d,"u_g"]&B\arrow[r]\arrow[d,equal]&0\\
0\arrow[r]&A'\arrow[r,"i_g"]&g_*E\arrow[r,"q_g"]&B\arrow[r]&0.
\end{tikzcd}
\]
The left square commutes because $[0,i(a)]=[g(a),0]$. The right square
commutes by the formula for $q_g$.

The bottom row is exact. If $[a',0]=0$, then
$(a',0)=(g(a),-i(a))$ for some $a$; injectivity of $i$ gives $a=0$
and then $a'=0$. If $q_g([a',e])=0$, write $e=i(a)$. The relation gives
\[
[a',e]=[a',i(a)]=[a'+g(a),0],
\]
so the kernel is the image of $i_g$. That image is visibly killed by
$q_g$, and $q_g$ is surjective because $q$ is. Again no injectivity or
surjectivity of $g$ was required.

The left square is a pushout. Given maps $r:A'\to T$ and $s:E\to T$
with $rg=si$, the formula
\[
[a',e]\longmapsto r(a')+s(e)
\]
defines a linear map $g_*E\to T$: the map on the direct sum kills
$(g(a),-i(a))$. It has the prescribed composites with $i_g$ and $u_g$.
Uniqueness follows because every class is $[a',0]+[0,e]$.

These two constructions have opposite directions for a reason. A map
\emph{into} the quotient permits us to restrict the extension over a new
domain. A map \emph{out of} the subobject permits us to impose a new
identification in the kernel. Their universal properties specify exactly
which compatible maps survive the construction.

\subsection{Adding extensions by merging their kernels}

Suppose
\[
\xi_j:\quad 0\longrightarrow A\xrightarrow{i_j}E_j
 \xrightarrow{q_j}B\longrightarrow0
\qquad (j=1,2)
\]
are two extensions with the same endpoints. Taking their direct sum gives
two copies of $B$ and two copies of $A$. We first require the quotient
coordinates to agree, and then merge the kernel coordinates by addition.
Thus we pull back along the diagonal $B\to B\oplus B$ and push out along
the sum $A\oplus A\to A$.

Here is the resulting object in one formula. Let
\[
T_2=\{(e_1,e_2):q_1(e_1)=q_2(e_2)\},\qquad
L_2=\{(i_1(a),-i_2(a)):a\in A\},
\]
and set $Q_2=T_2/L_2$. Both are submodules of the relevant direct sum,
and $L_2\subseteq T_2$. Define
\[
i_+(a)=[i_1(a),0]=[0,i_2(a)],\qquad
q_+([e_1,e_2])=q_1(e_1)=q_2(e_2).
\]
The equality defining $i_+$ is exactly the relation imposed by $L_2$.
The map $i_+$ is injective: if $(i_1(a),0)=(i_1(c),-i_2(c))$, then
$c=0$ and $a=0$. Every pair in the kernel of $q_+$ has the form
$(i_1(a),i_2(b))$, and its class is $i_+(a+b)$. Surjectivity follows by
lifting a given $b\in B$ through both $q_1$ and $q_2$. Hence
\[
\xi_1+\xi_2:\quad
0\longrightarrow A\xrightarrow{i_+}Q_2\xrightarrow{q_+}B
\longrightarrow0
\]
is an extension. This operation is called the \emph{Baer sum}.

To check the asserted pullback--pushout description, the pullback is
$T_2$, with kernel $A\oplus A$. Its pushout along addition is
\[
\frac{A\oplus T_2}
 {\{(a+b,(-i_1(a),-i_2(b))):a,b\in A\}}.
\]
It maps to $Q_2$ by
\[
[c,(e_1,e_2)]\longmapsto[e_1+i_1(c),e_2].
\]
A displayed relation maps to $[i_1(b),-i_2(b)]=0$. Its inverse sends
$[e_1,e_2]$ to $[0,(e_1,e_2)]$; the relation defining $L_2$ maps to a
pushout relation with $b=-a$. The two composites are identities, since
$[c,(e_1,e_2)]=[0,(e_1+i_1(c),e_2)]$ in the pushout. This verifies the
construction, including its endpoint maps.

\begin{proposition}\label{x:baer-group}
Baer sum makes the equivalence classes of extensions of $B$ by $A$ into
an abelian group. Its zero is the split extension. An inverse of
$(i,E,q)$ is $(-i,E,q)$.
\end{proposition}
\begin{proof}
First the operation is well-defined on classes. Equivalences
$h_j:E_j\to E'_j$ send $T_2$ into $T'_2$ and $L_2$ into $L'_2$
coordinatewise. They induce a map of extensions, which is an isomorphism
by Lemma~\ref{x:marked-iso}.

For associativity, form the common-image tuples for any $n\geq1$:
\[
T_n=\{(e_1,\ldots,e_n):q_1(e_1)=\cdots=q_n(e_n)\},
\]
and the relation submodule
\[
L_n=\{(i_1(a_1),\ldots,i_n(a_n)):a_1+\cdots+a_n=0\}.
\]
The quotient $Q_n=T_n/L_n$ is an extension of $B$ by $A$. Its quotient
map takes the common image; its kernel map sends $a$ to
$[i_1(a),0,\ldots,0]$. This kernel map is injective, because a tuple of
that form belongs to $L_n$ only if $a=0$. A tuple in the kernel of the
quotient map has entries $i_j(a_j)$; modulo $L_n$ it equals the kernel
image of $\sum_j a_j$. Finally, a given element of $B$ can be lifted in
each coordinate. These statements prove exactness.

In the following calculation, $E_1+E_2$ denotes the middle module of
the Baer sum, with its displayed endpoint maps.
For $n=3$ the map
\[
T_3\longrightarrow (E_1+E_2)+E_3,
\qquad (e_1,e_2,e_3)\longmapsto[[e_1,e_2],e_3]
\]
is linear. If $a_1+a_2+a_3=0$, its value on
$(i_1(a_1),i_2(a_2),i_3(a_3))$ is the class of
$(i_{12}(a_1+a_2),i_3(a_3))$, which is zero by the outer sum relation.
Here $i_{12}$ denotes the kernel inclusion for $E_1+E_2$.
Thus it factors through $Q_3$. The induced map is the identity on the
fixed endpoints, so it is an equivalence. The map
\[
(e_1,e_2,e_3)\longmapsto[e_1,[e_2,e_3]]
\]
likewise identifies $Q_3$ with $E_1+(E_2+E_3)$: a kernel tuple maps to
the relation determined by $a_1+(a_2+a_3)=0$. Both bracketings therefore
represent the same extension class. This proves associativity without
assuming it for the new operation.

Swapping the two coordinates identifies $E_1+E_2$ with $E_2+E_1$.
It preserves the kernel inclusion because
$[i_1(a),0]=[0,i_2(a)]$, and preserves the common quotient image.
This proves commutativity.

For the split extension $S=A\oplus B$, the formula
\[
E+S\longrightarrow E,\qquad [e,(a,b)]\longmapsto e+i(a)
\]
is defined on pairs with $q(e)=b$. The relation
$(i(c),(-c,0))$ maps to zero, and both endpoint maps are preserved.
Lemma~\ref{x:marked-iso} makes it an equivalence, proving the identity law.

Finally, in $E+(-E)$, where $-E$ denotes $(-i,E,q)$, the relation
submodule consists of pairs $(i(a),i(a))$. The linear map
$e\mapsto[e,e]$ kills $i(A)$, so it factors uniquely through $E/i(A)$.
The quotient map identifies $E/i(A)$ with $B$; hence we obtain a linear
map $s:B\to E+(-E)$ with $q_+s=\id_B$. By
Proposition~\ref{x:splitting}, this sum is split. The inverse law follows.
\end{proof}

These equivalence classes form a set. Indeed, choosing a set-theoretic
lift of every element of $B$ identifies the underlying set of a middle
module with $A\times B$, by $(a,b)\mapsto i(a)+s(b)$; this uses only
a choice of representatives, not a linear section. Module operations
and endpoint maps on this fixed set range over sets of functions.
Thus the extension classes can be represented by a set of such data.
For arbitrary modules this set-theoretic observation uses the usual
axiom of choice; none of the finite-filtration comparison proofs needs it.

We denote this abelian group by $\Ext^1_R(B,A)$. The order of the
arguments records the sequence: $A$ is the kernel and $B$ the quotient.
The superscript $1$ refers to a single extension step. No higher
extension theory is needed for what follows. These are classical
constructions; the notation and the general pullback--pushout description
are also recorded in the \href{https://stacks.math.columbia.edu/tag/010I}
{Stacks Project, Section 12.6}.

The new associativity has a concrete source. Three extensions can be
added by retaining triples with a common image and then imposing the
single condition that the three kernel coordinates sum to zero. Either
pairwise construction gives that same object up to the specified
isomorphism. The quotient has removed the dependence on how the addition
was parenthesized.

\subsection{An extension that fits in one matrix}\label{x:operator-extensions}

Let $k$ be a field. For $\lambda\in k$, let $S_\lambda$ be the
one-dimensional $k[t]$-module on which $t$ acts by multiplication by
$\lambda$. It is simple because a nonzero vector spans its underlying
one-dimensional space. Consider
\[
0\longrightarrow S_\mu\xrightarrow{i}E\xrightarrow{q}S_\lambda
\longrightarrow0.
\]
Write $T$ for the action of $t$. Choose $u=i(1)$ and choose $v\in E$
with $q(v)=1$. These form a $k$-basis: every $e\in E$ can first have its
quotient coordinate removed by subtracting $q(e)v$, leaving an element
of $i(S_\mu)$; and applying $q$ to $cu+dv=0$ gives $d=0$, followed by
$c=0$ since $i$ is injective.

The endpoint maps commute with the operators, so
\[
Tu=\mu u,\qquad Tv=\lambda v+a u
\]
for a unique $a\in k$. In this basis,
\[
[T]_{(u,v)}=\begin{pmatrix}\mu&a\\0&\lambda\end{pmatrix}.
\]
Conversely, for any $a$, these formulas define a linear operator on
$ku\oplus kv$, hence a $k[t]$-module $E_a$, and an extension with
$i(1)=u$, $q(u)=0$, and $q(v)=1$. Thus every extension under discussion
has this form.

The choice of $u$ is fixed by the marked endpoint, whereas another lift
of $1$ has the form $v'=v+bu$. Its coefficient is found by one calculation:
\[
T(v+bu)=\lambda(v+bu)+\bigl(a+b(\mu-\lambda)\bigr)u.
\]
The change of lift acts on the attachment coefficient by
\[
a\longmapsto a+b(\mu-\lambda).
\]

\begin{proposition}\label{x:operator-ext-classification}
The preceding coefficient gives an isomorphism of abelian groups
\[
\Ext^1_{k[t]}(S_\lambda,S_\mu)
 \cong k/(\mu-\lambda)k.
\]
In particular, extensions split when $\lambda\ne\mu$. When
$\lambda=\mu$, each $a\in k$ gives a distinct marked extension class,
and its class is zero exactly when $a=0$.
\end{proposition}
\begin{proof}
Changing the lift changes $a$ by an element of $(\mu-\lambda)k$, so
its coset is independent of that choice. An equivalence fixes $u$ and
sends a lift of $1$ to a lift of $1$; thus it preserves the same coset.
Conversely, if $a'=a+b(\mu-\lambda)$, the map $E_{a'}\to E_a$ sending
$u'\mapsto u$ and $v'\mapsto v+bu$ commutes with $T$ by the displayed
calculation. It therefore commutes with every polynomial in $T$, is
$k[t]$-linear, and is an equivalence of extensions. The classification
is consequently bijective.

To compute its addition, take $E_a$ and $E_c$ with chosen vectors
$u_1,v_1$ and $u_2,v_2$. In their Baer sum put
\[
u_+=[u_1,0]=[0,u_2],\qquad v_+=[v_1,v_2].
\]
The first vector is the marked kernel image of $1$ and the second maps
to $1$ in the quotient. The same spanning and independence argument
used above makes them a basis. Their operator satisfies
\[
Tu_+=\mu u_+,\qquad
Tv_+=\lambda v_++[a u_1,c u_2]
       =\lambda v_++(a+c)u_+.
\]
Thus Baer sum corresponds to addition of coefficient classes.

The split extension has coefficient zero. If $\mu-\lambda\ne0$, it
is invertible in $k$, so $(\mu-\lambda)k=k$ and the quotient group is
zero. If $\mu=\lambda$, the ideal is zero, giving the stated result.
\end{proof}

When the eigenvalues differ, the explicit section can be read off as
well. Choose $b=-a/(\mu-\lambda)$. Then $v'=v+bu$ satisfies
$Tv'=\lambda v'$, and the map $S_\lambda\to E$, $c\mapsto cv'$, is
$k[t]$-linear and is a section of $q$.

When the eigenvalues agree, a nonzero $a$ cannot be removed by changing
the lift. Yet all nonzero $a$ give isomorphic \emph{middle modules}.
Indeed, put $w=au$; then
\[
Tw=\lambda w,\qquad Tv=\lambda v+w,
\]
which has the same matrix for every nonzero $a$. This rescaling changes
the marked kernel vector, so it is not an equivalence keeping the
endpoints fixed. The $a=0$ middle module is different: $T-\lambda\id$
is zero there, whereas it is nonzero for $a\ne0$, a property preserved
by every intertwining isomorphism.

For $\lambda=\mu=1$, the cases $a=0$ and $a=1$ give the identity
operator and the triangular shear $(x,y)\mapsto(x+y,y)$. Both layers carry the
identity action. The coefficient $a$ now records explicitly the
attachment that their common list of layers forgot.

There are therefore two different classification questions. One asks
for the middle operator up to change of basis. The other remembers a
specified subobject and a specified quotient map. The latter contains
more information, which is why it supports the extension group just
constructed.

\subsection{A boundary map measures failure to lift}\label{x:obstruction}

Return to an arbitrary extension
\[
\xi:\quad0\longrightarrow A\xrightarrow{i}E\xrightarrow{q}B
\longrightarrow0.
\]
Given $f:P\to B$, we can now form a precise obstruction to finding
$\widetilde f:P\to E$ with $q\widetilde f=f$: take the pullback
extension and its class. Define
\[
\delta_\xi:\Hom_R(P,B)\longrightarrow\Ext^1_R(P,A),
\qquad f\longmapsto[f^*E].
\]
If $\widetilde f$ exists, then $p\mapsto(\widetilde f(p),p)$ is a
section of the pullback extension. Conversely, a section of its
projection to $P$ has exactly this form, and its first coordinate is
a lift. Proposition~\ref{x:splitting} therefore proves
\[
\delta_\xi(f)=0
\quad\Longleftrightarrow\quad
f\text{ has a linear lift through }q.
\]

This is a homomorphism of abelian groups. To verify the assertion,
take $f,g:P\to B$. The map
\[
(f^*E)+(g^*E)\longrightarrow(f+g)^*E,
\qquad[(e_1,p),(e_2,p)]\longmapsto(e_1+e_2,p)
\]
is well-defined: the target condition is
$q(e_1+e_2)=f(p)+g(p)$, and a Baer relation
$((i(a),0),(-i(a),0))$ maps to zero. It is linear and preserves both
endpoint maps. Lemma~\ref{x:marked-iso} makes it an equivalence, proving
$\delta_\xi(f+g)=\delta_\xi(f)+\delta_\xi(g)$. The pullback along
the zero map has section $p\mapsto(0,p)$, so $\delta_\xi(0)=0$.

In particular, we have proved exactness of the sequence of abelian groups
\[
0\longrightarrow\Hom_R(P,A)\xrightarrow{i\circ-}\Hom_R(P,E)
\xrightarrow{q\circ-}\Hom_R(P,B)
\xrightarrow{\delta_\xi}\Ext^1_R(P,A).
\]
Here exactness is asserted at the three $\Hom$ groups. The first map
is injective because $i$ is injective. A map $h:P\to E$ is killed by
$q\circ-$ exactly when its image lies in $i(A)$; it then factors
uniquely as $h=i\bar h$, with $\bar h$ linear because $i$ is an
injective linear map. Exactness at $\Hom_R(P,B)$ is the lifting
criterion proved above. No surjectivity onto $\Ext^1_R(P,A)$ is claimed.

For example, if $P=B$ and $f=\id_B$, its pullback is equivalent to
the original extension by $e\mapsto(e,q(e))$. Thus
$\delta_\xi(\id_B)=[\xi]$: lifting the identity is exactly splitting
the extension. In the equal-eigenvalue example, this class is the
coefficient $a$. A section would require an eigenvector above $1$;
the calculation shows exactly when such a vector exists.

There is also a useful contrast. If $P=R$, every map $f:R\to B$ lifts:
choose $e$ with $q(e)=f(1)$ and define $\widetilde f(r)=re$. This is
linear and satisfies $q\widetilde f(r)=r f(1)=f(r)$. The obstruction
depends on the object from which we test the quotient. A test object
with relations can detect constraints that a single free generator
cannot.

\subsection{An attachment can skip an adjacent layer}\label{x:skip-layer}

It is tempting to recover a filtered object by listing its simple
layers and the extension between every adjacent pair. Even that list
can miss information. The following example identifies the loss without
using any classification theorem.

On $k^3$ fix the flag
\[
0=F_0\subset F_1=ke_1\subset F_2=ke_1+ke_2\subset F_3=k^3.
\]
Compare the zero operator with the operator $N$ given by
\[
N(e_1)=N(e_2)=0,\qquad N(e_3)=e_1.
\]
Both preserve every $F_i$, so these are filtrations by $k[t]$-submodules.
Each of the three successive one-dimensional quotients has zero induced
operator in both examples, hence is $S_0$.

Now retain the adjacent two-layer pieces as well. On $F_2/F_0=F_2$
the operator is zero in both examples. On $F_3/F_1$ it is also zero,
because $N(F_3)\subseteq F_1$. Thus the marked extensions
\[
0\to F_1/F_0\to F_2/F_0\to F_2/F_1\to0,
\qquad
0\to F_2/F_1\to F_3/F_1\to F_3/F_2\to0
\]
are split and agree, using the indicated coordinate identifications,
for the two operators. Nevertheless the full modules are not
isomorphic. An isomorphism intertwining $N$ with zero would satisfy
$hN=0$; injectivity of $h$ would force $N=0$, contrary to $N(e_3)=e_1$.

The map from the third layer into the first is invisible in either
adjacent two-layer quotient. It remains visible if we retain the whole
extension
\[
0\longrightarrow F_2\longrightarrow F_3\longrightarrow F_3/F_2
\longrightarrow0
\]
with its maps: in the $N$ example every lift of the quotient generator
has the form $e_3+x$ with $x\in F_2$, and $N(e_3+x)=e_1\ne0$,
so no $k[t]$-linear section exists. For the zero operator the evident
section does exist.

The example specifies what the proposed summary discarded. Complete
successive extension data includes each already assembled module and
the maps attaching the next layer. Replacing those data by isolated
adjacent two-layer summaries can erase a longer attachment. We now
seek observations that retain enough information to reconstruct all
finite-dimensional nilpotent operators.
